% paper2_supplemental.tex
% Comprehensive Supplemental Material for:
% "Information Loss is Local:
%  The Gravitational Refractive Index as Petz Recovery Fidelity"
% Paper II in the tau framework series
% Author: Sheng-Kai Huang
% Compile: pdflatex paper2_supplemental.tex

\documentclass[aps,prd,preprint,superscriptaddress,longbibliography]{revtex4-2}

\usepackage{amsmath}
\usepackage{amssymb}
\usepackage{amsfonts}
\usepackage{amsthm}
\usepackage{bm}
\usepackage{hyperref}
\usepackage{booktabs}
\usepackage{graphicx}

\newtheorem{theorem}{Theorem}
\newtheorem{lemma}[theorem]{Lemma}
\newtheorem{proposition}[theorem]{Proposition}
\newtheorem{corollary}[theorem]{Corollary}
\newtheorem*{remark}{Remark}

% Convenience macros (same as main text)
\newcommand{\kB}{k_{\mathrm{B}}}
\newcommand{\Tr}{\mathrm{Tr}}
\newcommand{\id}{\mathrm{id}}
\newcommand{\Petz}{\mathcal{R}}
\newcommand{\chan}{\mathcal{N}}
\newcommand{\hilb}{\mathcal{H}}
\newcommand{\code}{\mathcal{C}}
\newcommand{\KL}{D}
\newcommand{\ket}[1]{|#1\rangle}
\newcommand{\bra}[1]{\langle#1|}
\newcommand{\braket}[2]{\langle#1|#2\rangle}
\newcommand{\op}[2]{|#1\rangle\langle#2|}
\newcommand{\abs}[1]{\left|#1\right|}
\newcommand{\norm}[1]{\left\|#1\right\|}
\newcommand{\tauasym}{\tau}

% Paper 2 specific macros
\newcommand{\ceff}{c_{\mathrm{eff}}}
\newcommand{\rs}{r_{\mathrm{s}}}
\newcommand{\Sgrav}{\Sigma_{\mathrm{grav}}}
\newcommand{\Fbound}{F_{\mathrm{bound}}}
\newcommand{\nref}{n_{\mathrm{grav}}}

\begin{document}

\title{Comprehensive Supplemental Material for\\
``Information Loss is Local:\\
The Gravitational Refractive Index as Petz Recovery Fidelity''}

\author{Sheng-Kai Huang}
\email{akai@fawstudio.com}
\affiliation{Independent Researcher}

\date{\today}

\maketitle

\tableofcontents

%=======================================================================
\section{S1. \texorpdfstring{$\Sgrav = -\ln(-g_{00})$}{Sigma\_grav = -ln(-g00)} as a Theorem}
\label{sec:three_routes}
%=======================================================================

A hostile reading of our main identification
$\ceff/c = e^{-\Sgrav/2}$ might dismiss it as a tautology:
``you \emph{defined} $\Sgrav$ to make this true.''
This section shows that $\Sgrav = -\ln(-g_{00})$ is a
\emph{theorem}---derivable from standard results in quantum information
theory and general relativity---not a definition.

The argument has three independent legs:
\begin{itemize}
\item \textbf{Route C (Channel Theorem):} The gravitational redshift
  \emph{is} a pure-loss bosonic channel with transmissivity
  $\eta = -g_{00}$.  The unique extensive entropy production rate
  of such a channel is $-\ln\eta$.
  This is a theorem in quantum optics combined with standard GR.
\item \textbf{Route B (Landauer):} The Tolman relation forces
  the dimensionless erasure cost ratio to be $(-g_{00})^{-1/2}$;
  the logarithm of the cost ratio is $-\ln(-g_{00})$.
\item \textbf{Route A (Modular flow):} The entanglement first law
  gives the exact QRE ratio $D_{\mathrm{out}}/D_{\mathrm{in}}
  = -g_{00}$; the extensive measure $-\ln(D_{\mathrm{out}}/D_{\mathrm{in}})$
  then equals $-\ln(-g_{00})$.
\end{itemize}
All three yield $\Sgrav = -\ln(-g_{00})$.  We present Route~C first
because it provides the cleanest theorem, then Routes~A and~B as
independent confirmations.

%-----------------------------------------------------------------------
\subsection{Route C (Channel Theorem): Gravitational redshift as
pure-loss channel}
\label{sec:routeC}
%-----------------------------------------------------------------------

This is the route that most directly answers the tautology
objection.  It proceeds in five steps; every step is either a
theorem in quantum information theory (Steps~1--3) or a standard
result in general relativity (Steps~4--5).

\emph{Step 1 (QI---no gravity): Pure-loss bosonic channel.}---A
single-mode pure-loss bosonic channel with transmissivity
$\eta \in (0,1]$ is the CPTP map~\cite{Weedbrook2012}
\begin{equation}
  \chan_\eta(\rho) = \sum_{k=0}^{\infty} E_k\,\rho\,E_k^\dagger\,,
  \quad
  E_k = \sqrt{\binom{n}{k}}\;\eta^{n/2}(1-\eta)^{k/2}\;\ket{n-k}\bra{n}\,.
  \label{eq:Kraus_loss}
\end{equation}
This is the standard quantum-optical model of a beam
splitter~\cite{Weedbrook2012}.

\emph{Step 2 (QI---no gravity): QRE under pure loss.}---For coherent
states $\ket{\alpha}$ and $\ket{\beta}$, the quantum relative entropy
is $D(\ket{\alpha}\|\ket{\beta}) = |\alpha - \beta|^2$.
After the channel:
\begin{equation}
  D(\chan_\eta\ket{\alpha}\|\chan_\eta\ket{\beta})
  = \eta\,|\alpha - \beta|^2
  = \eta\,D_{\mathrm{in}}\,.
  \label{eq:DeltaD_loss}
\end{equation}
This is exact (not perturbative).  The QRE ratio is
\begin{equation}
  \frac{D_{\mathrm{out}}}{D_{\mathrm{in}}} = \eta\,.
  \label{eq:QRE_ratio}
\end{equation}

\emph{Step 3 (QI---no gravity): Extensivity forces the
logarithm.}---Two sequential pure-loss channels with
transmissivities $\eta_1$ and $\eta_2$ compose to a single
pure-loss channel with transmissivity $\eta_1\eta_2$.
Define the \emph{entropy production} of a channel as a functional
$\Sigma[\chan]$ satisfying:
\begin{enumerate}
  \item[(i)] $\Sigma[\chan] \geq 0$, with equality iff
    $\chan$ is an isometry;
  \item[(ii)] $\Sigma[\chan_2\circ\chan_1]
    = \Sigma[\chan_1] + \Sigma[\chan_2]$
    (extensivity / additivity under composition).
\end{enumerate}
\begin{proposition}[Uniqueness of extensive entropy production]
\label{prop:extensivity}
The unique functional on pure-loss channels satisfying (i)--(ii) is
\begin{equation}
  \Sigma[\chan_\eta] = -\ln\eta\,,
  \label{eq:Sigma_unique}
\end{equation}
up to a positive multiplicative constant (set to~1 by
normalization to nats).
\end{proposition}
\begin{proof}
Property~(ii) requires $\Sigma[\chan_{\eta_1\eta_2}]
= \Sigma[\chan_{\eta_1}] + \Sigma[\chan_{\eta_2}]$.
Writing $f(\eta) \equiv \Sigma[\chan_\eta]$, this is
$f(\eta_1\eta_2) = f(\eta_1) + f(\eta_2)$, the Cauchy
functional equation on $(0,1]$.  With the monotonicity
constraint $f'(\eta) \leq 0$ (from (i) and the data
processing inequality: smaller $\eta$ means more loss,
hence more entropy production), the unique continuous
solution is $f(\eta) = -c\ln\eta$ with $c > 0$.
\end{proof}

\begin{remark}
The na\"ive fractional loss $\Delta D/D_{\mathrm{in}} = 1 - \eta$
is \emph{not} extensive: for two channels,
$1 - \eta_1\eta_2 \neq (1-\eta_1) + (1-\eta_2)$.
Extensivity uniquely selects the logarithm.
\end{remark}

\emph{Step 4 (GR---no QI): Gravitational frequency shift.}---A photon
of frequency $\omega_r$ emitted at radius $r$ in a static spacetime
is detected at infinity with frequency
\begin{equation}
  \omega_\infty = \sqrt{-g_{00}(r)}\;\omega_r\,.
  \label{eq:freq_shift}
\end{equation}
This is a standard result of general relativity, exact for any
static metric~\cite{MTW1973}.

\emph{Step 5 (GR---no QI): Transmissivity equals $-g_{00}$.}---The
gravitational redshift maps a single-mode coherent state
$\ket{\alpha}$ to $\ket{\sqrt{\eta}\,\alpha}$ with
\begin{equation}
  \eta = \frac{\omega_\infty^2}{\omega_r^2} = -g_{00}(r)\,,
  \label{eq:eta}
\end{equation}
where the $\omega^2$ factor accounts for both the frequency shift and
the time dilation of the emission rate (energy $\propto \omega$,
flux $\propto \omega$)~\cite{AhmadiFuentes2014}.
This identifies the gravitational redshift as a pure-loss
bosonic channel $\chan_\eta$ with $\eta = -g_{00}$.

\begin{theorem}[Channel Theorem]
\label{thm:channel}
Let $(M,g)$ be a static spacetime with time-time metric
component $g_{00}(r)$.  The gravitational redshift from
radius $r$ to spatial infinity defines a pure-loss bosonic
channel $\chan_{-g_{00}}$.  Its unique extensive entropy
production (Proposition~\ref{prop:extensivity}) is
\begin{equation}
  \boxed{\Sgrav(r) = -\ln(-g_{00}(r))\,.}
  \label{eq:channel_theorem}
\end{equation}
\end{theorem}
\begin{proof}
Combine Steps~1--5.  Steps~1--3 are theorems in quantum
information theory (no reference to gravity).  Steps~4--5
are standard general relativity (no reference to QI).
The only ``identification'' is recognizing that the
gravitational redshift is a pure-loss channel, which
follows from the Kraus decomposition being the standard
beam-splitter with $\eta = -g_{00}$.
\end{proof}

\begin{remark}[Why this is not a tautology]
We did not \emph{define} $\Sgrav \equiv -\ln(-g_{00})$.
Rather: (a)~GR dictates $\eta = -g_{00}$; (b)~quantum
information theory dictates that the unique extensive
entropy production of a pure-loss channel with
transmissivity $\eta$ is $-\ln\eta$; (c)~combining
(a) and (b) \emph{forces} $\Sgrav = -\ln(-g_{00})$.
The logarithm is not a choice---it is the unique
answer demanded by extensivity.
\end{remark}

\emph{Consistency check: channel composition.}---A photon
propagating from $r_1$ to $r_2 > r_1$ and then from $r_2$
to infinity traverses two sequential channels with
transmissivities $\eta_1 = -g_{00}(r_1)/(-g_{00}(r_2))$ and
$\eta_2 = -g_{00}(r_2)$.  The composed transmissivity is
$\eta_1 \eta_2 = -g_{00}(r_1)$, and the entropy productions
add:
\begin{equation}
  \Sgrav(r_1 \to \infty) = \Sgrav(r_1 \to r_2) + \Sgrav(r_2 \to \infty)\,,
  \label{eq:composition}
\end{equation}
i.e., $-\ln(-g_{00}(r_1)) = [-\ln(-g_{00}(r_1)) + \ln(-g_{00}(r_2))]
+ [-\ln(-g_{00}(r_2))]$.  This is trivially satisfied, confirming
extensivity.  The fractional loss $1-\eta$ does \emph{not} satisfy
this: $(1-(-g_{00}(r_1))) \neq
[1-(-g_{00}(r_1))/(-g_{00}(r_2))] + [1-(-g_{00}(r_2))]$.

\emph{Remark on the factor of 2.}---In the amplitude convention
$\sqrt{\eta_{\mathrm{amp}}} = \omega_\infty/\omega_r = \sqrt{-g_{00}}$,
the transmissivity is $\eta_{\mathrm{amp}} = -g_{00}$ and the entropy
production is $-\ln(\eta_{\mathrm{amp}}) = -\ln(-g_{00})$.  In the
frequency-only convention
$\eta_{\mathrm{freq}} = (-g_{00})^{1/2}$ (without the time-dilation
contribution), one would get
$\Sigma = -\frac{1}{2}\ln(-g_{00})$.  The correct answer requires
including both the frequency shift and the time dilation of the
emission rate, giving $\eta = -g_{00}$ and $\Sigma = -\ln(-g_{00})$.

%-----------------------------------------------------------------------
\subsection{Route A: Modular flow and entanglement first law}
\label{sec:routeA}
%-----------------------------------------------------------------------

This route uses the entanglement first law in the AQFT setting
to derive the QRE ratio between two radii, and then invokes
the extensivity argument of Proposition~\ref{prop:extensivity}
to obtain $\Sgrav$.  Steps~1--4 below use standard results from
AQFT~\cite{BisognanoWichmann1976} and the Tolman
relation~\cite{Tolman1930}; the modular Hamiltonian structure
is closely related to the framework of Dorau and
Much~\cite{DorauMuch2025}.

\emph{Step 1: Bisognano--Wichmann theorem.}---For a Rindler wedge in
Minkowski spacetime, the modular operator $\Delta_\Omega$ associated
with the vacuum state $\ket{\Omega}$ and the wedge algebra
$\mathcal{A}(W)$ generates Lorentz boosts~\cite{BisognanoWichmann1976}:
\begin{equation}
  \Delta_\Omega^{is} = e^{-2\pi K s}\,,
  \label{eq:BW}
\end{equation}
where $K$ is the boost generator and $s$ is the modular flow parameter.

\emph{Step 2: Extension to Schwarzschild.}---For a static observer at
radius $r$ in the Schwarzschild spacetime, the local Rindler
approximation gives the modular Hamiltonian
\begin{equation}
  H_{\mathrm{mod}}(r) = \frac{2\pi}{\kappa\sqrt{-g_{00}(r)}}\,H_{\mathrm{local}}\,,
  \label{eq:Hmod}
\end{equation}
where $\kappa$ is the surface gravity and the factor
$1/\sqrt{-g_{00}(r)}$ is the Tolman blueshift (using
the general form valid for any static metric, not only
Schwarzschild).

\emph{Step 3: Entanglement first law.}---For small perturbations
$\delta\rho$ about the vacuum, the QRE equals the modular energy:
\begin{equation}
  D(\rho\|\sigma_\beta)
  = \Tr[\delta\rho\,H_{\mathrm{mod}}]
  = \frac{2\pi}{\kappa\sqrt{-g_{00}(r)}}\,\langle\delta E\rangle_r\,.
  \label{eq:first_law}
\end{equation}

\emph{Step 4: QRE ratio.}---An excitation of local energy
$\langle\delta E\rangle_r$ at radius $r$ is observed at infinity as
$\langle\delta E\rangle_\infty = \sqrt{-g_{00}(r)}\;
\langle\delta E\rangle_r$ (gravitational redshift).
Substituting:
\begin{equation}
  D_{\mathrm{in}}(r)
  = \frac{2\pi}{\kappa\,(-g_{00}(r))}\,\langle\delta E\rangle_\infty\,,
  \qquad
  D_{\mathrm{out}}(\infty) = \frac{2\pi}{\kappa}\,\langle\delta E\rangle_\infty\,.
  \label{eq:Din_Dout}
\end{equation}
The QRE ratio is therefore
\begin{equation}
  \frac{D_{\mathrm{out}}}{D_{\mathrm{in}}} = -g_{00}(r)\,.
  \label{eq:QRE_ratio_A}
\end{equation}
This is \emph{exact} for any static metric---not only Schwarzschild.

\emph{Step 5: From ratio to entropy production.}---The QRE ratio
$D_{\mathrm{out}}/D_{\mathrm{in}} = -g_{00} = \eta$ is identical
to the channel transmissivity of Route~C
(Eq.~\eqref{eq:QRE_ratio}).  By the extensivity argument
(Proposition~\ref{prop:extensivity}):
\begin{equation}
  \boxed{\Sgrav^{(\mathrm{a})} = -\ln\!\left(\frac{D_{\mathrm{out}}}{D_{\mathrm{in}}}\right)
  = -\ln(-g_{00})\,.}
  \label{eq:route_a_result}
\end{equation}

\begin{remark}[Fractional loss vs.\ logarithmic loss]
The \emph{fractional} QRE loss is
$\Delta D/D_{\mathrm{in}} = 1 - (-g_{00})$, which for
Schwarzschild gives $\rs/r$.  The \emph{logarithmic} entropy
production is $-\ln(-g_{00})$, which for Schwarzschild gives
$-\ln(1 - \rs/r) = \rs/r + \tfrac{1}{2}(\rs/r)^2 + \cdots$.
These agree at first order in $\rs/r$ but differ at
$O((\rs/r)^2)$.  The logarithmic form is uniquely selected by
extensivity (additivity under channel composition).  An earlier
version of this derivation incorrectly presented the boxed
result as following directly from the fractional loss;
we thank the referee for pressing this point.
\end{remark}

%-----------------------------------------------------------------------
\subsection{Route B: Gravitational Landauer principle (Herrera 2020)}
\label{sec:routeB}
%-----------------------------------------------------------------------

This route uses the thermodynamic cost of information erasure in
a gravitational field.

\emph{Step 1: Landauer principle at infinity.}---At spatial infinity
(flat space), the minimum energy cost to erase one bit of information
at temperature $T_\infty$ is~\cite{Landauer1961}
\begin{equation}
  W_\infty = \kB T_\infty \ln 2\,.
  \label{eq:Landauer_inf}
\end{equation}

\emph{Step 2: Tolman relation.}---In thermal equilibrium in a static
gravitational field, the local temperature satisfies the Tolman
relation~\cite{Tolman1930}:
\begin{equation}
  T(r)\,\sqrt{-g_{00}(r)} = T_\infty\,.
  \label{eq:Tolman}
\end{equation}
Therefore
\begin{equation}
  T(r) = \frac{T_\infty}{\sqrt{-g_{00}(r)}}\,.
  \label{eq:Tr}
\end{equation}

\emph{Step 3: Landauer cost at radius $r$.}---The minimum erasure cost
at radius $r$ is
\begin{equation}
  W(r) = \kB T(r)\ln 2 = \frac{\kB T_\infty \ln 2}{\sqrt{-g_{00}(r)}}
  = \frac{W_\infty}{\sqrt{-g_{00}(r)}}\,.
  \label{eq:Wr}
\end{equation}
The excess cost relative to infinity is
\begin{equation}
  \Delta W = W(r) - W_\infty
  = W_\infty\!\left(\frac{1}{\sqrt{-g_{00}}} - 1\right).
  \label{eq:DeltaW}
\end{equation}

\emph{Step 4: Dimensionless entropy production.}---Following
Herrera~\cite{Herrera2020}, the dimensionless gravitational entropy
production is twice the logarithm of the cost ratio (factor of 2
because both the forward and reverse Landauer costs contribute):
\begin{equation}
  \Sgrav^{(\mathrm{b})} = 2\ln\!\left(\frac{W(r)}{W_\infty}\right)
  = 2\ln\!\left(\frac{1}{\sqrt{-g_{00}}}\right)
  = -\ln(-g_{00})\,.
  \label{eq:Sigma_b}
\end{equation}

\emph{Result.}---
\begin{equation}
  \boxed{\Sgrav^{(\mathrm{b})} = -\ln(-g_{00})\,.}
  \label{eq:route_b_result}
\end{equation}
For the exponential metric: $\Sgrav = \rs/r$.  For Schwarzschild:
$\Sgrav = -\ln(1-\rs/r) \approx \rs/r + \frac{1}{2}(\rs/r)^2 + \cdots$.

\emph{Physical interpretation.}---The gravitational entropy production
is the excess informational cost of performing a bit-erasure operation
at radius $r$ compared to infinity.  Gravity makes information
processing more expensive at lower radii because the local temperature
is higher (Tolman relation), and erasing a bit against a hotter bath
requires more energy.

%-----------------------------------------------------------------------
\subsection{Convergence of the three routes}
\label{sec:convergence}
%-----------------------------------------------------------------------

All three routes yield the same result, but with different logical
structures:

\begin{table}[h]
\caption{\label{tab:routes}Summary of the three derivation routes.
The ``Theorem?''\ column indicates whether the route is a
self-contained theorem (T) or uses an identification step (I).}
\begin{ruledtabular}
\begin{tabular}{lllll}
  \textbf{Route} & \textbf{Physics} & \textbf{Key input}
    & \textbf{Result} & \textbf{T/I} \\
  \colrule
  C (channel) & QI + GR redshift
    & $\eta = -g_{00}$, extensivity & $-\ln(-g_{00})$ & T \\
  A (modular) & AQFT + entanglement 1st law
    & $D_{\mathrm{out}}/D_{\mathrm{in}} = -g_{00}$ & $-\ln(-g_{00})$ & T \\
  B (Landauer) & Tolman + erasure cost
    & $W(r)/W_\infty = 1/\sqrt{-g_{00}}$ & $-\ln(-g_{00})$ & T \\
\end{tabular}
\end{ruledtabular}
\end{table}

The common mathematical structure is:
\begin{enumerate}
  \item Each route independently derives a \emph{ratio} involving
  $-g_{00}$: the channel transmissivity (C), the QRE ratio (A),
  or the Landauer cost ratio (B).
  \item The logarithm of each ratio gives $-\ln(-g_{00})$.  In
  Routes~A and~C, the logarithm is forced by extensivity
  (Proposition~\ref{prop:extensivity}).  In Route~B, it arises
  because the Landauer cost is proportional to temperature,
  and the Tolman relation gives $T(r)/T_\infty = (-g_{00})^{-1/2}$,
  so the entropy production $= 2\ln(T(r)/T_\infty) = -\ln(-g_{00})$.
\end{enumerate}

The convergence of three routes---quantum-optical (channel),
algebraic (modular flow), and thermodynamic (Landauer)---with
the same result $\Sgrav = -\ln(-g_{00})$ provides strong evidence
that this is not an artifact of any single framework.  Route~C
(Theorem~\ref{thm:channel}) is the cleanest, as every step is
either a standard theorem in quantum information theory or a
standard result in general relativity, and the combination is
a theorem, not an identification.

%=======================================================================
\section{S2. PPN Expansion of the Exponential Metric}
\label{sec:PPN}
%=======================================================================

We provide a complete parametrized post-Newtonian (PPN) analysis of the
exponential metric, demonstrating that it passes all current
solar-system tests.

%-----------------------------------------------------------------------
\subsection{PPN parameters}
%-----------------------------------------------------------------------

In isotropic coordinates, the exponential metric is
\begin{align}
  g_{00}^{\mathrm{exp}} &= -e^{-\rs/r}
  = -\!\left(1 - \frac{\rs}{r} + \frac{1}{2}\frac{\rs^2}{r^2}
    - \frac{1}{6}\frac{\rs^3}{r^3} + \cdots\right),
  \label{eq:g00_exp}\\
  g_{ij}^{\mathrm{exp}} &= e^{+\rs/r}\,\delta_{ij}
  = \left(1 + \frac{\rs}{r} + \frac{1}{2}\frac{\rs^2}{r^2}
    + \cdots\right)\delta_{ij}\,.
  \label{eq:gij_exp}
\end{align}
Writing $U = \rs/(2r) = GM/(c^2 r)$, the standard PPN form is
\begin{align}
  g_{00} &= -(1 - 2U + 2\beta U^2 + \cdots)\,,
  \label{eq:PPN_g00}\\
  g_{ij} &= (1 + 2\gamma U + \cdots)\,\delta_{ij}\,.
  \label{eq:PPN_gij}
\end{align}
Comparing Eq.~\eqref{eq:g00_exp} with Eq.~\eqref{eq:PPN_g00}:
\begin{equation}
  -g_{00} = 1 - 2U + 2U^2 + \cdots
  \quad\Longrightarrow\quad \beta = 1\,.
  \label{eq:beta}
\end{equation}
Comparing Eq.~\eqref{eq:gij_exp} with Eq.~\eqref{eq:PPN_gij}:
\begin{equation}
  g_{ij} = (1 + 2U + \cdots)\,\delta_{ij}
  \quad\Longrightarrow\quad \gamma = 1\,.
  \label{eq:gamma}
\end{equation}

\begin{equation}
  \boxed{\gamma = \beta = 1 \quad\text{(identical to GR at 1PN).}}
  \label{eq:PPN_result}
\end{equation}

%-----------------------------------------------------------------------
\subsection{2PN differences}
%-----------------------------------------------------------------------

At second post-Newtonian (2PN) order, the exponential and
Schwarzschild metrics differ.  Schwarzschild in isotropic coordinates:
\begin{align}
  g_{00}^{\mathrm{Schw}} &= -\frac{(1-U)^2}{(1+U)^2}
  = -(1 - 4U + 8U^2 - 16U^3 + \cdots)\,,
  \label{eq:g00_Schw_iso}\\
  g_{ij}^{\mathrm{Schw}} &= (1+U)^4\,\delta_{ij}
  = (1 + 4U + 6U^2 + \cdots)\,\delta_{ij}\,.
  \label{eq:gij_Schw_iso}
\end{align}
Note that in standard isotropic form, Schwarzschild uses
$U_{\mathrm{iso}} = M/(2r_{\mathrm{iso}})$, and the expansion
coefficients differ from our $U = \rs/(2r)$ convention.

In the convention $U = GM/(c^2 r)$ (where $r$ is the isotropic radial
coordinate), the expansions are:
\begin{center}
\begin{tabular}{lcc}
\toprule
\textbf{Coefficient} & \textbf{Exponential} & \textbf{Schwarzschild} \\
\midrule
$g_{00}$: $U^1$ & $-2$ & $-2$ \\
$g_{00}$: $U^2$ & $+2$ & $+2$ \\
$g_{00}$: $U^3$ & $-\frac{4}{3}$ & $-\frac{8}{3}$ \\
$g_{ij}$: $U^1$ & $+2$ & $+2$ \\
$g_{ij}$: $U^2$ & $+2$ & $+\frac{3}{2}$ \\
\bottomrule
\end{tabular}
\end{center}
The first difference is at $O(U^2)$ in $g_{ij}$ and $O(U^3)$ in $g_{00}$.

%-----------------------------------------------------------------------
\subsection{Mercury precession}
%-----------------------------------------------------------------------

The perihelion precession per orbit for a general static spherically
symmetric metric with PPN parameters $\gamma$ and $\beta$ is
\begin{equation}
  \Delta\phi = \frac{6\pi GM}{c^2 a(1-e^2)}
  \left(\frac{2-\beta+2\gamma}{3}\right)
  + O(U^2)\,,
  \label{eq:precession}
\end{equation}
where $a$ is the semi-major axis and $e$ the eccentricity.  For both
Schwarzschild and the exponential metric, $\gamma = \beta = 1$, giving
\begin{equation}
  \Delta\phi_{\mathrm{exp}} = \Delta\phi_{\mathrm{Schw}}
  = \frac{6\pi GM}{c^2 a(1-e^2)} = 42.98''/\mathrm{cy}\,.
  \label{eq:Mercury}
\end{equation}

At 2PN, the correction to Mercury's precession is~\cite{Mychelkin2024}:
\begin{equation}
  \delta(\Delta\phi) = \Delta\phi_{\mathrm{exp}} - \Delta\phi_{\mathrm{Schw}}
  \sim \frac{6\pi GM}{c^2 a(1-e^2)}\cdot O\!\left(\frac{GM}{c^2 a}\right)
  \approx 2.3\times 10^{-7}\;\mathrm{arcsec/cy}\,.
  \label{eq:Mercury_2PN}
\end{equation}
Current measurement precision is
$\sim 3 \times 10^{-3}$~arcsec/cy~\cite{Bertotti2003}, so the 2PN
difference is four orders of magnitude below detectability.

%-----------------------------------------------------------------------
\subsection{Light bending: debunking the ``half bending'' myth}
\label{sec:light_bending}
%-----------------------------------------------------------------------

A persistent misconception holds that the exponential metric
``predicts half the observed light bending.''  This conflates the
exponential metric (a solution of GR + phantom scalar) with the
Yilmaz field equations (an alternative theory of gravity).

The deflection angle for light in a general static, spherically
symmetric metric with PPN parameters $\gamma$ and $\beta$ is
\begin{equation}
  \delta\theta = \frac{(1+\gamma)}{2}\cdot\frac{4GM}{c^2 b}
  + O(U^2)\,,
  \label{eq:bending}
\end{equation}
where $b$ is the impact parameter.  Since $\gamma = 1$ for the
exponential metric (proven above):
\begin{equation}
  \delta\theta_{\mathrm{exp}} = \frac{4GM}{c^2 b}
  = \delta\theta_{\mathrm{Schw}}\,.
  \label{eq:bending_full}
\end{equation}

The ``half bending'' confusion arises from evaluating the bending
contribution from $g_{00}$ alone, which gives $2GM/(c^2 b)$.  In GR,
an equal contribution comes from the spatial curvature ($g_{ij}$
term).  The exponential metric has $\gamma = 1$, so both contributions
are present and equal:
\begin{equation}
  \delta\theta = \underbrace{\frac{2GM}{c^2 b}}_{g_{00}\;\text{contribution}}
  + \underbrace{\frac{2GM}{c^2 b}}_{g_{ij}\;\text{contribution}}
  = \frac{4GM}{c^2 b}\,.
  \label{eq:bending_decomposed}
\end{equation}

The Cassini bound $|\gamma - 1| < 2.3 \times 10^{-5}$~\cite{Bertotti2003}
is satisfied exactly by the exponential metric.

%=======================================================================
\section{S3. Echo Delay Calculation}
\label{sec:echo_detail}
%=======================================================================

We compute the gravitational-wave echo delay for the exponential
metric in detail.

%-----------------------------------------------------------------------
\subsection{Tortoise coordinate}
%-----------------------------------------------------------------------

The tortoise coordinate for the exponential metric
$g_{00} = -e^{-\rs/r}$, $g_{rr} = e^{+\rs/r}$ (isotropic
coordinates) is
\begin{equation}
  r^* = \int \frac{\sqrt{g_{rr}}}{\sqrt{-g_{00}}}\,dr
  = \int e^{\rs/r}\,dr\,.
  \label{eq:tortoise_exp}
\end{equation}
This integral is finite for all $r > 0$, in sharp contrast to the
Schwarzschild tortoise coordinate
$r^*_{\mathrm{Schw}} = r + \rs\ln|r/\rs - 1| \to -\infty$ as
$r \to \rs$.

%-----------------------------------------------------------------------
\subsection{The echo integral}
%-----------------------------------------------------------------------

The round-trip light travel time from the photon sphere at
$r_{\mathrm{ph}} = \rs$ (isotropic coordinates) to the wormhole
throat at $r_{\mathrm{throat}} = \rs/2$ and back is
\begin{equation}
  \Delta t_{\mathrm{echo}} = \frac{2}{c}\int_{r_{\mathrm{throat}}}^{r_{\mathrm{ph}}}
  e^{\rs/r}\,dr
  = \frac{2}{c}\int_{\rs/2}^{\rs} e^{\rs/r}\,dr\,.
  \label{eq:echo_integral}
\end{equation}
Substituting $u = \rs/r$ (so $r = \rs/u$, $dr = -\rs\,du/u^2$):
\begin{equation}
  \Delta t_{\mathrm{echo}} = \frac{2\rs}{c}\int_{1}^{2}
  \frac{e^u}{u^2}\,du\,.
  \label{eq:echo_u}
\end{equation}

\emph{Numerical evaluation.}---The integral can be evaluated using
the exponential integral function $\mathrm{Ei}(x)$:
\begin{equation}
  \int \frac{e^u}{u^2}\,du = -\frac{e^u}{u} + \mathrm{Ei}(u) + C\,.
  \label{eq:Ei}
\end{equation}
Therefore
\begin{align}
  \int_1^2 \frac{e^u}{u^2}\,du
  &= \left[-\frac{e^u}{u} + \mathrm{Ei}(u)\right]_1^2
  \nonumber\\
  &= \left(-\frac{e^2}{2} + \mathrm{Ei}(2)\right)
  - \left(-e + \mathrm{Ei}(1)\right)
  \nonumber\\
  &= -\frac{e^2}{2} + e + \mathrm{Ei}(2) - \mathrm{Ei}(1)\,.
  \label{eq:echo_exact}
\end{align}
Using $e^2/2 \approx 3.695$, $e \approx 2.718$,
$\mathrm{Ei}(2) \approx 4.954$, $\mathrm{Ei}(1) \approx 1.895$:
\begin{equation}
  \int_1^2 \frac{e^u}{u^2}\,du \approx -3.695 + 2.718 + 4.954 - 1.895
  = 2.082\,.
  \label{eq:echo_numerical}
\end{equation}
The echo delay is
\begin{equation}
  \boxed{\Delta t_{\mathrm{echo}} \approx \frac{4.16\,\rs}{c}\,.}
  \label{eq:echo_result}
\end{equation}

%-----------------------------------------------------------------------
\subsection{Numerical estimates for astrophysical sources}
%-----------------------------------------------------------------------

\begin{table}[h]
\caption{\label{tab:echo_detail}Echo delays for representative sources.
$\rs = 2GM/c^2$.}
\begin{ruledtabular}
\begin{tabular}{lccc}
  \textbf{Source} & $M/M_\odot$ & $\rs$ (km)
    & $\Delta t_{\mathrm{echo}}$ \\
  \colrule
  GW150914 remnant & 62 & 183
    & 2.5~ms \\
  GW190521 remnant & 150 & 443
    & 6.1~ms \\
  Intermediate BH & $10^3$ & 2,950
    & 41~ms \\
  Sgr~A* & $4\times10^6$ & $1.18\times10^7$
    & 163~s \\
  M87* & $6.5\times10^9$ & $1.92\times10^{10}$
    & $7.4\times10^4$~s ($\approx 21$~hr)\\
\end{tabular}
\end{ruledtabular}
\end{table}

%-----------------------------------------------------------------------
\subsection{Comparison with Planck-wall echo model}
%-----------------------------------------------------------------------

In the ``Planck wall'' model of Cardoso et~al.~\cite{Cardoso2016},
the reflecting surface is placed at Planck-length distance from
the would-be horizon.  The echo delay in Schwarzschild tortoise
coordinates is
\begin{equation}
  \Delta t_{\mathrm{Planck}} \approx \frac{\rs}{c}
  \ln\!\left(\frac{\rs}{\ell_P}\right),
  \label{eq:echo_Planck}
\end{equation}
where $\ell_P \approx 1.6\times 10^{-35}$~m is the Planck length.
For a $60\,M_\odot$ remnant:
$\ln(\rs/\ell_P) \approx \ln(10^{38}) \approx 88$, giving
$\Delta t_{\mathrm{Planck}} \approx 88\,\rs/c \approx 50$~ms.

The exponential-metric echo is approximately $20\times$ shorter
than the Planck-wall echo:
\begin{equation}
  \frac{\Delta t_{\mathrm{echo}}}{\Delta t_{\mathrm{Planck}}}
  \approx \frac{4.2}{\ln(\rs/\ell_P)} \sim \frac{1}{20}\,.
  \label{eq:ratio_echo}
\end{equation}
This difference is observationally significant: current echo
searches~\cite{Abedi2017} use templates with $\Delta t \sim 200$~ms,
appropriate for Planck-wall models but potentially missing the
shorter exponential-metric signal.

%=======================================================================
\section{S4. Petz Saturation Analysis}
\label{sec:saturation}
%=======================================================================

%-----------------------------------------------------------------------
\subsection{Why exact saturation is impossible for $\Sigma > 0$}
%-----------------------------------------------------------------------

\begin{theorem}[No exact saturation for $\Sigma > 0$]
\label{thm:no_saturation}
Let $\chan$ be a CPTP map and $\sigma$ a faithful state.  The JRSWW
bound~\cite{JRSWW2018}
\begin{equation}
  F(\rho,\,\tilde{\Petz}_{\sigma,\chan}\circ\chan(\rho))^2
  \geq \exp(-\Delta D)
  \label{eq:JRSWW}
\end{equation}
is saturated (equality holds for all $\rho$) if and only if
$\Delta D = 0$, i.e., $\chan$ is sufficient for $\sigma$.
\end{theorem}

\begin{proof}
The JRSWW proof proceeds via the Golden--Thompson
inequality~\cite{GoldenThompson1965}:
\begin{equation}
  \Tr[e^{A+B}] \leq \Tr[e^A e^B]\,,
  \label{eq:GT}
\end{equation}
with equality if and only if $[A,B] = 0$.  In the application to the
Petz recovery bound, $A$ and $B$ involve $\log\sigma$ and
$\log\chan(\sigma)$; these commute if and only if the channel is
sufficient for $\sigma$, which by Petz's theorem occurs precisely when
$\Delta D = 0$~\cite{Petz1988}.

More precisely, the chain of inequalities in the JRSWW proof is:
\begin{align}
  &F(\rho,\tilde{\Petz}\circ\chan(\rho))^2
  \nonumber\\
  &\quad\geq \exp\!\left(-D_M(\rho\|\sigma)\right)
  \cdot \exp\!\left(D_M(\chan(\rho)\|\chan(\sigma))\right)
  \nonumber\\
  &\quad= \exp(-\Delta D_M) \geq \exp(-\Delta D)\,,
  \label{eq:chain}
\end{align}
where $D_M$ is the measured relative entropy and
$\Delta D_M \leq \Delta D$ with equality for commuting states.
The first inequality uses the Golden--Thompson inequality and
becomes an equality only when $[\log\sigma,\,
\chan^\dagger(\log\chan(\sigma))] = 0$, which for faithful $\sigma$
is equivalent to $\Delta D = 0$.
\end{proof}

\begin{remark}
This theorem means that for any gravitational channel with
$\Sgrav > 0$, the Petz recovery fidelity strictly exceeds
$\exp(-\Sgrav/2)$.  The exponential metric, defined by promoting
$F = \exp(-\Sgrav/2)$ to an identity, should be understood as a
\emph{lower bound geometry}---the metric that would result if
recovery were as poor as the JRSWW bound allows.  Nature may do
better, and the actual metric may lie between the exponential and
Schwarzschild predictions.
\end{remark}

%-----------------------------------------------------------------------
\subsection{Gap calculation}
%-----------------------------------------------------------------------

For a weak gravitational field ($\Sgrav = \rs/r \ll 1$), we can
estimate the gap between the actual Petz recovery fidelity and the
JRSWW bound.

Consider the amplitude damping channel as a model for gravitational
loss.  The amplitude damping channel with parameter
$\gamma = 1 - \eta$ (where $\eta = e^{-\rs/r} \approx 1 - \rs/r$
for weak fields) has Kraus operators
\begin{equation}
  K_0 = \begin{pmatrix} 1 & 0 \\ 0 & \sqrt{\eta} \end{pmatrix},
  \qquad
  K_1 = \begin{pmatrix} 0 & \sqrt{1-\eta} \\ 0 & 0 \end{pmatrix}.
  \label{eq:AD_Kraus}
\end{equation}

For the maximally mixed reference state $\sigma = I/2$, the Petz
recovery map is
\begin{equation}
  \Petz_{\sigma,\chan}(\rho) = \frac{1}{2}\sum_k K_k^\dagger
  \left(\frac{I}{2}\right)^{-1/2}\rho
  \left(\frac{I}{2}\right)^{-1/2} K_k
  = \sum_k K_k^\dagger \rho\, K_k\,.
  \label{eq:Petz_AD}
\end{equation}
The fidelity $F_{\mathrm{Petz}}$ for the input $\ket{+} = (\ket{0}+\ket{1})/\sqrt{2}$
can be computed explicitly.  Writing $\epsilon = 1-\eta = \rs/r$:
\begin{align}
  F_{\mathrm{Petz}}^2 &= 1 - \frac{1}{4}\epsilon^2 + O(\epsilon^3)\,,
  \label{eq:F_Petz_AD}\\
  \exp(-\Sgrav) &= \exp(-\epsilon) = 1 - \epsilon + \frac{1}{2}\epsilon^2 + O(\epsilon^3)\,.
  \label{eq:exp_Sgrav}
\end{align}
The gap is
\begin{equation}
  \boxed{\Delta \equiv F_{\mathrm{Petz}}^2 - \exp(-\Sgrav)
  = \epsilon - \frac{3}{4}\epsilon^2 + O(\epsilon^3)
  = \frac{\rs}{r} - \frac{3}{4}\left(\frac{\rs}{r}\right)^2 + \cdots\,.}
  \label{eq:gap}
\end{equation}
The leading gap is $O(\rs/r)$, not $O((\rs/r)^2)$.  This means the
Petz recovery does significantly better than the JRSWW bound in weak
fields.  The relative gap
$\Delta/\exp(-\Sgrav) \approx \rs/r$ is small, confirming
approximate saturation.

%-----------------------------------------------------------------------
\subsection{Schwarzschild fidelity vs.\ exponential fidelity}
%-----------------------------------------------------------------------

For the Schwarzschild metric, $\Fbound = \sqrt{1-\rs/r}$:
\begin{equation}
  F_{\mathrm{Schw}} = \sqrt{1-\rs/r}
  = 1 - \frac{1}{2}\frac{\rs}{r} - \frac{1}{8}\left(\frac{\rs}{r}\right)^2 + \cdots
  \label{eq:F_Schw}
\end{equation}
For the exponential metric, $\Fbound = e^{-\rs/(2r)}$:
\begin{equation}
  F_{\mathrm{exp}} = e^{-\rs/(2r)}
  = 1 - \frac{1}{2}\frac{\rs}{r} + \frac{1}{8}\left(\frac{\rs}{r}\right)^2 + \cdots
  \label{eq:F_exp}
\end{equation}
The difference is
\begin{equation}
  F_{\mathrm{Schw}} - F_{\mathrm{exp}}
  = -\frac{1}{4}\left(\frac{\rs}{r}\right)^2 + O\!\left(\left(\frac{\rs}{r}\right)^3\right).
  \label{eq:F_diff}
\end{equation}
At the neutron star surface ($\rs/r \approx 0.4$):
$|F_{\mathrm{Schw}} - F_{\mathrm{exp}}| \approx 0.04$, a $5\%$
difference.  At $\rs/r = 1$ (horizon):
$F_{\mathrm{Schw}} = 0$ but $F_{\mathrm{exp}} = e^{-1/2} \approx 0.607$.

%=======================================================================
\section{S5. Vaidya Calculation: \texorpdfstring{$\tauasym = 1$}{tau = 1} at the Apparent Horizon}
\label{sec:vaidya}
%=======================================================================

We provide the complete derivation of the Kodama vector and its
connection to $\tauasym$ in the Vaidya spacetime, complementing the
abbreviated treatment in the existing supplement.

%-----------------------------------------------------------------------
\subsection{Vaidya metric setup}
%-----------------------------------------------------------------------

The ingoing Vaidya metric in advanced Eddington--Finkelstein
coordinates $(v,r,\theta,\phi)$ is
\begin{equation}
  ds^2 = -f(v,r)\,dv^2 + 2\,dv\,dr + r^2\,d\Omega^2\,,
  \label{eq:vaidya_full}
\end{equation}
where $f(v,r) = 1 - 2M(v)/r$ and $M(v)$ is the Bondi mass.

The areal radius is $R = r$.  The 2D orbit-space metric on $(v,r)$ is
\begin{equation}
  h_{AB}\,dx^A\,dx^B = -f\,dv^2 + 2\,dv\,dr\,.
  \label{eq:h2d_full}
\end{equation}

The Kodama vector is $K_A = \epsilon_A{}^B\,\partial_B R$ with
$\epsilon_{vr} = 1$, giving (as derived in the companion supplement):
\begin{equation}
  K^a\partial_a = -\partial_v\,,\qquad |K| = \sqrt{f} = \sqrt{1 - 2M(v)/r}\,.
  \label{eq:Kodama_result}
\end{equation}

%-----------------------------------------------------------------------
\subsection{$g_{00} \to 0$ at the MOTS}
%-----------------------------------------------------------------------

The marginally outer trapped surface (MOTS) is defined by the
vanishing of the expansion of outgoing null geodesics.  For the Vaidya
metric~\eqref{eq:vaidya_full}, the outgoing null expansion is
\begin{equation}
  \theta_+ = \frac{2}{r}\left(1 - \frac{2M(v)}{r}\right)^{1/2}
  \cdot\frac{1}{r}\left(r - 2M(v)\right)\,.
  \label{eq:theta_plus}
\end{equation}
More precisely, the expansion of the outgoing null normal
$\ell^a = \partial_v + \frac{f}{2}\partial_r$ is
\begin{equation}
  \theta_+ = \frac{f}{r} = \frac{1}{r}\left(1 - \frac{2M(v)}{r}\right).
  \label{eq:theta_plus_simple}
\end{equation}
This vanishes when $r = 2M(v)$, i.e., at the MOTS (apparent horizon).

At the MOTS:
\begin{equation}
  f\big|_{r=2M(v)} = 1 - \frac{2M(v)}{2M(v)} = 0\,,
  \label{eq:f_MOTS}
\end{equation}
so $|K| = 0$ and
\begin{equation}
  \tauasym_K = 1 - |K| = 1 - 0 = 1\,.
  \label{eq:tau_MOTS}
\end{equation}
Complete information loss at the apparent horizon.

%-----------------------------------------------------------------------
\subsection{Connection to quantum extremal surface}
%-----------------------------------------------------------------------

The quantum extremal surface (QES)~\cite{Engelhardt2015} extremizes
the generalized entropy $S_{\mathrm{gen}} = A/(4G\hbar) + S_{\mathrm{bulk}}$.
In the semiclassical limit, the QES coincides with the MOTS.
Including quantum corrections:
\begin{equation}
  r_{\mathrm{QES}} = 2M(v) + \frac{4G\hbar}{r}\cdot
  \frac{\partial S_{\mathrm{bulk}}}{\partial r}\bigg|_{r=2M(v)}
  + O(\hbar^2)\,.
  \label{eq:QES_correction}
\end{equation}
The quantum-corrected $\tauasym$ at the QES is
\begin{equation}
  \tauasym_K(r_{\mathrm{QES}}) = 1 - O(\hbar/M^2) < 1\,,
  \label{eq:tau_QES}
\end{equation}
so the quantum correction makes the information loss slightly
incomplete---consistent with the expectation that quantum gravity
effects prevent complete information loss.

%=======================================================================
\section{S6. Pikovski Channel Analysis}
\label{sec:pikovski}
%=======================================================================

We provide a complete worked example of the Pikovski gravitational
dephasing channel, including the Petz recovery map, fidelity
calculations, and a detailed argument for why it does \emph{not}
yield the gravitational channel $\chan_{\mathrm{grav}}$.

%-----------------------------------------------------------------------
\subsection{The Pikovski dephasing channel}
%-----------------------------------------------------------------------

Pikovski et~al.~\cite{Pikovski2015} showed that gravitational time
dilation universally dephases composite quantum systems.  A particle
in a superposition of heights $\Delta h$ in a gravitational field
$g$ experiences a phase difference
\begin{equation}
  \Delta\phi = \frac{m c^2\,g\,\Delta h\,t}{\hbar c^2}
  = \frac{\Delta E\,\Delta\Phi\,t}{\hbar c^2}\,,
  \label{eq:Pikovski_phase}
\end{equation}
where $\Delta E$ is the internal energy spread and
$\Delta\Phi = g\,\Delta h$ is the gravitational potential difference.

For a system with internal energy distribution $\{E_n, p_n\}$, the
coherence between height states $\ket{h_1}$ and $\ket{h_2}$ decays as
\begin{equation}
  \rho_{12}(t) = \rho_{12}(0)\cdot\chi(\Delta t_{\mathrm{grav}})\,,
  \label{eq:coherence_decay}
\end{equation}
where $\chi(\tau) = \sum_n p_n\,e^{-i E_n \tau/\hbar}$ is the
characteristic function of the internal energy distribution and
$\Delta t_{\mathrm{grav}} = g\,\Delta h\,t/c^2$ is the
proper-time difference.

\emph{Coherence parameter.}---The dephasing channel has a single
coherence parameter
\begin{equation}
  p = |\chi(\Delta t_{\mathrm{grav}})|\,.
  \label{eq:p_coherence}
\end{equation}

\emph{Two-level system at zero temperature.}---For a two-level
internal system with gap $\omega_0$:
\begin{equation}
  p = |\cos(\omega_0\,\Delta t_{\mathrm{grav}}/2)|\,.
  \label{eq:p_2level}
\end{equation}

\emph{Thermal state.}---For a thermal internal state at temperature $T$:
\begin{equation}
  p = \left|\frac{Z(\beta - i\Delta t_{\mathrm{grav}}/\hbar)}{Z(\beta)}\right|,
  \label{eq:p_thermal}
\end{equation}
where $Z(\beta) = \Tr[e^{-\beta H}]$ is the partition function.

%-----------------------------------------------------------------------
\subsection{Kraus operators}
%-----------------------------------------------------------------------

The qubit dephasing channel with parameter $p$ has Kraus operators
\begin{equation}
  K_0 = \sqrt{\frac{1+p}{2}}\,I\,,\qquad
  K_1 = \sqrt{\frac{1-p}{2}}\,\sigma_z\,,
  \label{eq:dephasing_Kraus}
\end{equation}
satisfying $K_0^\dagger K_0 + K_1^\dagger K_1 = I$.  The channel
acts as
\begin{equation}
  \chan_p(\rho) = p\,\rho + (1-p)\,\frac{\sigma_z\rho\sigma_z + \rho}{2}
  = \begin{pmatrix} \rho_{00} & p\,\rho_{01} \\ p\,\rho_{10} & \rho_{11}
  \end{pmatrix}.
  \label{eq:dephasing_action}
\end{equation}

%-----------------------------------------------------------------------
\subsection{Petz recovery map for dephasing}
%-----------------------------------------------------------------------

For the maximally mixed reference $\sigma = I/2$:

\emph{Step 1:} $\chan_p(\sigma) = \sigma = I/2$.

\emph{Step 2:} $\chan_p(\sigma)^{-1/2} = \sqrt{2}\,I$.

\emph{Step 3:} The Petz map is
\begin{align}
  \Petz_{\sigma,\chan_p}(\rho) &= \sigma^{1/2}\,\chan_p^\dagger
  \!\left(\chan_p(\sigma)^{-1/2}\,\rho\,\chan_p(\sigma)^{-1/2}\right)
  \sigma^{1/2}
  \nonumber\\
  &= \frac{1}{2}\,\chan_p^\dagger(2\rho)
  = \chan_p(\rho)\,.
  \label{eq:Petz_dephasing}
\end{align}
The key observation: for diagonal $\sigma$, the Petz recovery map
for the dephasing channel \emph{equals the channel itself}.  The
Petz map applies dephasing \emph{again}, making things worse rather
than better.

%-----------------------------------------------------------------------
\subsection{Fidelity calculations}
%-----------------------------------------------------------------------

For the input state $\ket{+} = (\ket{0}+\ket{1})/\sqrt{2}$:

\emph{After the channel:}
$\chan_p(\ket{+}\!\bra{+}) = \frac{1}{2}\left(\begin{smallmatrix}
1 & p \\ p & 1\end{smallmatrix}\right)$.

\emph{After Petz recovery:}
\begin{equation}
  \Petz\circ\chan_p(\ket{+}\!\bra{+})
  = \chan_p^2(\ket{+}\!\bra{+})
  = \frac{1}{2}\begin{pmatrix} 1 & p^2 \\ p^2 & 1 \end{pmatrix}.
  \label{eq:Petz_twice}
\end{equation}

\emph{Petz fidelity:}
\begin{equation}
  F_{\mathrm{Petz}} = \bra{+}\!\Petz\circ\chan_p(\ket{+}\!\bra{+})\!\ket{+}
  = \frac{1+p^2}{2}\,.
  \label{eq:F_Petz_dephasing}
\end{equation}

\emph{Optimal fidelity (identity map---do nothing):}
\begin{equation}
  F_{\mathrm{opt}} = \bra{+}\!\chan_p(\ket{+}\!\bra{+})\!\ket{+}
  = \frac{1+p}{2}\,.
  \label{eq:F_opt_dephasing}
\end{equation}

The Petz map is \emph{worse} than doing nothing:
$F_{\mathrm{Petz}} < F_{\mathrm{opt}}$ for $0 < p < 1$.  This is
because the Petz map is designed for the reference state $\sigma$,
not for the input state $\rho$; for dephasing with diagonal $\sigma$,
the Petz map coincides with the noise itself.

\emph{The JRSWW bound:}
\begin{equation}
  \Sigma_p = D(\ket{+}\!\bra{+}\|\sigma)
  - D(\chan_p(\ket{+}\!\bra{+})\|\chan_p(\sigma))
  = \ln 2 - H_{\mathrm{bin}}\!\left(\frac{1+p}{2}\right),
  \label{eq:Sigma_dephasing}
\end{equation}
where $H_{\mathrm{bin}}$ is the binary entropy.  For $p$ close to 1
($\epsilon = 1-p \ll 1$):
$\Sigma_p \approx \epsilon^2/(2\ln 2)$.

%-----------------------------------------------------------------------
\subsection{Why Pikovski is NOT the gravitational channel}
%-----------------------------------------------------------------------

There are three fundamental mismatches between the Pikovski dephasing
channel and the gravitational entropy production $\Sgrav = \rs/r$:

\emph{Mismatch 1: Probe dependence.}---The Pikovski entropy production
\begin{equation}
  \Sigma_{\mathrm{Pik}} \propto (\Delta E)^2\,(\Delta\Phi)^2\,t^2
  \label{eq:Sigma_Pik}
\end{equation}
depends on the internal energy spread $\Delta E$, the mass $m$,
and the observation time $t$.  In contrast, $\Sgrav = \rs/r$
is \emph{universal}---independent of the probe.

\emph{Mismatch 2: Divergent supremum.}---Taking
$\Delta E \to \infty$ gives $\Sigma_{\mathrm{Pik}} \to \infty$
at any fixed $\Delta\Phi$ and $t$.  No finite supremum exists.
In contrast, $\Sgrav$ is finite for all $r > 0$.

\emph{Mismatch 3: Basis.}---The Pikovski channel dephases in the
\emph{position} basis of center-of-mass motion.  The gravitational
channel should dephase (or lose information) in the
\emph{energy/frequency} basis, corresponding to the redshift of
photon frequencies.

\emph{What the Pikovski channel IS.}---The Pikovski dephasing rate per
unit frequency per unit time is
$\Gamma = g\,\Delta h/c^2 = \Delta\Phi/c^2 \sim \Sgrav/2$.  This
universal rate is correctly predicted by $\Sgrav$, but the
\emph{accumulated} dephasing depends on the specific probe.  The
Pikovski channel is a \emph{consequence} of gravitational information
loss, not the gravitational channel itself.

%-----------------------------------------------------------------------
\subsection{Numerical estimates}
%-----------------------------------------------------------------------

For a molecule of mass $m = 10^4$~amu with internal energy spread
$\Delta E = 1$~eV at height difference $\Delta h = 1$~m on Earth's
surface ($g \approx 9.8$~m/s$^2$):

\begin{equation}
  \Delta t_{\mathrm{grav}} = \frac{g\,\Delta h}{c^2}\,t
  \approx 1.1\times10^{-16}\,t\;\text{(seconds)}\,,
  \label{eq:dt_numerical}
\end{equation}
\begin{equation}
  p = \cos(\omega_0\,\Delta t_{\mathrm{grav}}/2)
  \approx 1 - \frac{(\Delta E)^2(\Delta t_{\mathrm{grav}})^2}{8\hbar^2}\,.
  \label{eq:p_numerical}
\end{equation}

For $t = 1$~s: $\Delta t_{\mathrm{grav}} \approx 10^{-16}$~s,
$\omega_0 \approx 1.5\times10^{15}$~rad/s (for $\Delta E = 1$~eV),
$\omega_0\,\Delta t_{\mathrm{grav}}/2 \approx 7.6\times10^{-2}$~rad,
$p \approx 0.997$, $\tauasym_{\mathrm{Pik}} \approx 0.003$.

The corresponding gravitational $\tauasym$ at the Earth's surface is
$\tauasym_{\mathrm{grav}} = 1 - \sqrt{1-r_{\mathrm{s},\oplus}/R_\oplus}
\approx 7\times10^{-10}$---six orders of magnitude smaller.  This
confirms that the Pikovski dephasing accumulates much faster than the
static gravitational $\tauasym$ for a probe-dependent reason (the
internal energy spread amplifies the effect).

%=======================================================================
\section{S7. Raychaudhuri Decomposition}
\label{sec:raychaudhuri}
%=======================================================================

%-----------------------------------------------------------------------
\subsection{Raychaudhuri equation and entropy production}
%-----------------------------------------------------------------------

For a congruence of timelike geodesics with tangent vector $u^a$,
the Raychaudhuri equation reads~\cite{Raychaudhuri1955,Wald1984}
\begin{equation}
  \frac{d\theta}{d\tau} = -\frac{\theta^2}{3} - \sigma_{ab}\sigma^{ab}
  + \omega_{ab}\omega^{ab} - R_{ab}\,u^a u^b\,,
  \label{eq:Raychaudhuri}
\end{equation}
where $\theta$ is the expansion scalar, $\sigma_{ab}$ the shear
tensor, and $\omega_{ab}$ the vorticity tensor.

Interpreting each term as a contribution to the rate of change of
gravitational entropy production~\cite{Paper2}:
\begin{equation}
  \frac{d\Sgrav}{ds} = \underbrace{\frac{\theta^2}{3}}_{\text{expansion}}
  + \underbrace{\sigma_{ab}\sigma^{ab}}_{\text{shear}}
  - \underbrace{\omega_{ab}\omega^{ab}}_{\text{vorticity}}
  + \underbrace{R_{ab}\,u^a u^b}_{\text{curvature}}\,.
  \label{eq:dSigma_ds}
\end{equation}

%-----------------------------------------------------------------------
\subsection{Static vacuum Schwarzschild: $\Sigma_{\mathrm{grav,Raych}} = 0$}
%-----------------------------------------------------------------------

For static observers in the Schwarzschild vacuum:
\begin{itemize}
\item $\theta = 0$ (static congruence, no expansion);
\item $\sigma_{ab} = 0$ (spherical symmetry, no shear);
\item $\omega_{ab} = 0$ (hypersurface-orthogonal, no vorticity);
\item $R_{ab} = 0$ (vacuum solution).
\end{itemize}
Therefore $d\Sgrav/ds = 0$: the Raychaudhuri decomposition gives
\emph{zero} entropy production rate.

This does \emph{not} contradict $\Sgrav = \rs/r$; rather, it shows
that $\Sgrav$ from the modular flow route is an \emph{integrated}
quantity (the total QRE decrease from $r$ to $\infty$), while the
Raychaudhuri decomposition gives the \emph{local rate}.  The two are
complementary, not contradictory:
\begin{equation}
  \Sgrav = \int_r^\infty \frac{d\Sgrav}{ds}\,ds + \text{boundary term}\,.
  \label{eq:integrated_vs_local}
\end{equation}
The boundary term (from the Tolman redshift at infinity) provides
$\Sgrav = -\ln(-g_{00})$; the local Raychaudhuri rate vanishes in
vacuum.

%-----------------------------------------------------------------------
\subsection{NEC violation implies $\tauasym < 0$}
%-----------------------------------------------------------------------

\begin{proposition}[One-directional NEC--$\tauasym$ connection]
If the null energy condition (NEC) is violated along a null geodesic
congruence, i.e., $R_{ab}\,k^a k^b < 0$ for some null vector $k^a$,
then there exist configurations where $d\Sgrav/d\lambda < 0$, implying
$\tauasym$ can decrease along the congruence.
\end{proposition}

\begin{proof}
For a null geodesic congruence with tangent $k^a$, the null
Raychaudhuri equation (with vorticity vanishing by the Frobenius
theorem for hypersurface-orthogonal null congruences) is
\begin{equation}
  \frac{d\theta}{d\lambda} = -\frac{\theta^2}{2}
  - \hat{\sigma}_{ab}\hat{\sigma}^{ab} - R_{ab}\,k^a k^b\,.
  \label{eq:null_Raychaudhuri}
\end{equation}
If $R_{ab}\,k^a k^b < 0$ (NEC violation), the last term is positive,
potentially making $d\theta/d\lambda > 0$ (defocusing).

In the entropy production interpretation:
\begin{equation}
  \frac{d\Sgrav}{d\lambda} = \frac{\theta^2}{2}
  + \hat{\sigma}_{ab}\hat{\sigma}^{ab} + R_{ab}\,k^a k^b\,.
  \label{eq:null_dSigma}
\end{equation}
If NEC is violated strongly enough that
$|R_{ab}\,k^a k^b| > \theta^2/2 + \hat{\sigma}^2$, then
$d\Sgrav/d\lambda < 0$: entropy production decreases, and $\tauasym$
can decrease.  This corresponds to information flowing \emph{back}
from the environment to the system---a form of non-Markovian
backflow.
\end{proof}

\begin{remark}
The converse does not hold: $d\Sgrav/d\lambda < 0$ does not require
NEC violation (it could result from decreasing shear or expansion).
The NEC--$\tauasym$ connection is one-directional:
NEC violation $\Rightarrow$ possible $\tauasym$ decrease, but not
conversely.
\end{remark}

%-----------------------------------------------------------------------
\subsection{Kerr vorticity contribution}
%-----------------------------------------------------------------------

For the Kerr metric with spin parameter $a$, the congruence of
static observers at Boyer--Lindquist radius $r$ has nonzero vorticity:
\begin{equation}
  \omega^2 = \omega_{ab}\omega^{ab}
  = \frac{2M^2 a^2 r^2}{\Sigma^4}\,,
  \label{eq:Kerr_vorticity}
\end{equation}
where $\Sigma^2 = r^2 + a^2\cos^2\theta$.  The vorticity term
$-\omega^2$ in the Raychaudhuri equation acts to \emph{decrease}
$d\Sgrav/ds$:
\begin{equation}
  \frac{d\Sgrav}{ds}\bigg|_{\mathrm{Kerr}} = R_{ab}\,u^a u^b - \omega^2
  < R_{ab}\,u^a u^b\,.
  \label{eq:Kerr_dSigma}
\end{equation}
Frame dragging (vorticity) reduces the effective entropy production
rate, consistent with the physical intuition that angular momentum
protects information.  In the extremal Kerr limit ($a \to M$), the
vorticity contribution becomes significant near the (would-be)
horizon, potentially allowing $\tauasym < 1$ at $r = r_+$.

%=======================================================================
% BIBLIOGRAPHY
%=======================================================================

\begin{thebibliography}{40}

\bibitem{BisognanoWichmann1976}
J.~J.~Bisognano and E.~H.~Wichmann,
``On the Duality Condition for Quantum Fields,''
J.\ Math.\ Phys.\ \textbf{17}, 303 (1976).

\bibitem{DorauMuch2025}
L.~Dorau and A.~Much,
``From Quantum Relative Entropy to the Semiclassical Einstein Equations,''
Phys.\ Rev.\ Lett.\ \textbf{135}, 241401 (2025); arXiv:2510.24491.

\bibitem{Landauer1961}
R.~Landauer,
``Irreversibility and Heat Generation in the Computing Process,''
IBM J.\ Res.\ Dev.\ \textbf{5}, 183 (1961).

\bibitem{Tolman1930}
R.~C.~Tolman,
``On the Weight of Heat and Thermal Equilibrium in General Relativity,''
Phys.\ Rev.\ \textbf{35}, 904 (1930).

\bibitem{Herrera2020}
L.~Herrera,
``Landauer Principle and General Relativity,''
Entropy \textbf{22}, 340 (2020).

\bibitem{AhmadiFuentes2014}
M.~Ahmadi, D.~E.~Bruschi, C.~Sabin, G.~Adesso, and I.~Fuentes,
``Relativistic Quantum Metrology,''
Sci.\ Rep.\ \textbf{4}, 4996 (2014).

\bibitem{MTW1973}
C.~W.~Misner, K.~S.~Thorne, and J.~A.~Wheeler,
\emph{Gravitation} (W.~H.~Freeman, San Francisco, 1973).

\bibitem{Weedbrook2012}
C.~Weedbrook, S.~Pirandola, R.~Garc\'ia-Patr\'on, N.~J.~Cerf,
T.~C.~Ralph, J.~H.~Shapiro, and S.~Lloyd,
``Gaussian quantum information,''
Rev.\ Mod.\ Phys.\ \textbf{84}, 621 (2012).

\bibitem{Mychelkin2024}
E.~G.~Mychelkin, M.~A.~Makukov, and G.~Suliyeva,
``On the weak and strong field effects in antiscalar background,''
Gen.\ Relativ.\ Gravit.\ \textbf{56}, 44 (2024); arXiv:2403.11610.

\bibitem{Bertotti2003}
B.~Bertotti, L.~Iess, and P.~Tortora,
``A test of general relativity using radio links with the Cassini
spacecraft,''
Nature \textbf{425}, 374 (2003).

\bibitem{Cardoso2016}
V.~Cardoso, E.~Franzato, and P.~Pani,
``Is the Gravitational-Wave Ringdown a Probe of the Event Horizon?,''
Phys.\ Rev.\ Lett.\ \textbf{116}, 171101 (2016).

\bibitem{Abedi2017}
J.~Abedi, H.~Dykaar, and N.~Afshordi,
``Echoes from the Abyss: Tentative evidence for Planck-scale structure
at black hole horizons,''
Phys.\ Rev.\ D \textbf{96}, 082004 (2017).

\bibitem{JRSWW2018}
M.~Junge, R.~Renner, D.~Sutter, M.~M.~Wilde, and A.~Winter,
``Universal Recovery Maps and Approximate Sufficiency of Quantum
Relative Entropy,''
Ann.\ Henri Poincar\'e \textbf{19}, 2955 (2018).

\bibitem{GoldenThompson1965}
S.~Golden, ``Lower bounds for the Helmholtz function,''
Phys.\ Rev.\ \textbf{137}, B1127 (1965);
C.~J.~Thompson, ``Inequality with applications in statistical
mechanics,'' J.\ Math.\ Phys.\ \textbf{6}, 1812 (1965).

\bibitem{Petz1988}
D.~Petz,
``Sufficiency of channels over von Neumann algebras,''
Q.\ J.\ Math.\ \textbf{39}, 97 (1988).

\bibitem{Engelhardt2015}
N.~Engelhardt and A.~C.~Wall,
``Quantum extremal surfaces: holographic entanglement entropy beyond
the classical regime,''
J.\ High Energy Phys.\ \textbf{2015}(01), 073 (2015).

\bibitem{Pikovski2015}
I.~Pikovski, M.~Zych, F.~Costa, and \v{C}.~Brukner,
``Universal decoherence due to gravitational time dilation,''
Nat.\ Phys.\ \textbf{11}, 668 (2015).

\bibitem{Paper2}
S.-K. Huang,
``Information Loss is Local: The Gravitational Refractive Index as Petz
Recovery Fidelity,'' companion paper (2026).

\bibitem{Raychaudhuri1955}
A.~Raychaudhuri,
``Relativistic Cosmology.\ I,''
Phys.\ Rev.\ \textbf{98}, 1123 (1955).

\bibitem{Wald1984}
R.~M.~Wald,
\emph{General Relativity} (University of Chicago Press, 1984).

\end{thebibliography}

\end{document}
